\subsection{Unresolved intervals and incumbent quality}
\label{sec:optimality-gaps}
The OPT/FEASIBLE distinction describes the outcome of a particular run;
it does not, by itself, determine the quality of a FEASIBLE incumbent.
We therefore conduct a separate post-hoc analysis of all
\GapObservations{} confirmation observations. For instance $i$, let
$\ell_i^{\mathrm{th}}$ be its applicable theoretical lower bound,
$\ell_{ir}$ the recorded lower bound in observation $r$, and $u_{ir}$
the independently validated span of that observation's labeling. Combining
the six observations gives
\[
 L_i=\max\bigl\{\ell_i^{\mathrm{th}},\max_r\ell_{ir}\bigr\},
 \qquad U_i=\min_r u_{ir},
 \qquad L_i\leq\lambda_{h,k}(G_i)\leq U_i.
\]
The lower-bound inequality relies on the recorded solver bounds where
they exceed the theoretical bound; this analysis does not independently
verify UNSAT or MILP optimality certificates. The interval combines
evidence across backends and repetitions and is not attributed to a single
30-second run. The original statuses and timed coverage are unchanged.

Of the \GapInstances{} instances, \GapClosed{} have closed intervals:
\GapTheory{} attain the applicable theoretical lower bound, and
\GapRecorded{} additional instances have a recorded OPT result.
No additional interval closes solely by combining bounds from FEASIBLE
runs. The remaining \GapOpen{} instances all belong to the sampled ER
groups; their interval widths range from \GapOpenMin{} to \GapOpenMax{}
label units. This identifies unresolved cases in this cohort, not a general
hardness ranking of graph families. A wide interval may reflect weak
lower bounds, poor incumbents, or both.

\begin{table}[htbp]
\centering\small
\caption{Post-hoc interval classification on the full confirmation cohort.
Theory denotes attainment of an applicable theoretical lower bound;
Other closed denotes additional closed intervals. These columns and Open
partition Cases. They do not count successes by either backend within
one run.}
\label{tab:gap-coverage}
\input{\GapReportDir/summary.tex}
\end{table}

For any stored witness of span $u$, the remaining suboptimality satisfies
\[
 u-U_i\leq u-\lambda_{h,k}(G_i)\leq u-L_i.
\]
If $L_i=U_i=u$, the witness is retrospectively known to be optimal even
when its original run returned FEASIBLE. If $u>U_i$, a strictly better
validated labeling exists, so the witness is known to be suboptimal even
if the optimum remains unresolved. Otherwise its optimality is unknown.
These are absolute bounds in label units, not confidence intervals or
the relative MIP gaps reported by a commercial solver.

\begin{table}[htbp]
\centering\small
\caption{Retrospective quality of observations originally reported as
FEASIBLE. Optimal witness attains a closed combined interval; Worse witness
has a strictly better stored witness; Unresolved means neither conclusion
is established. Repetitions are observations of the same instances.}
\label{tab:feasible-quality}
\input{\GapReportDir/quality.tex}
\end{table}

In particular, \GapSATFeasibleOptimal{} SAT FEASIBLE observations contain
optimal witnesses according to the combined evidence. They remain FEASIBLE
in every original coverage and runtime table: finding an optimal labeling
and proving its optimality within a run are distinct outcomes. The analysis
does not establish why a run failed to prove optimality or isolate the cost
of proving lower bounds. Such causal conclusions require further controlled
experiments. All unresolved instances, bound sources, and observation-level
classifications are exported with source hashes under
\path{results/analysis/confirmation_gaps/}.
\FloatBarrier
